% Lösungen Einheit 3 von 3 – Argumente über Formel und Verträglichkeit: die Länge des Intervalls bei doppeltem Umfang (Faktor eins durch Wurzel zwei (zusammenbau v0.1)
\einheitenkopf{Einheit 3 von 3 · Argumente über Formel und Verträglichkeit: die Länge des Intervalls bei doppeltem Umfang (Faktor eins durch Wurzel zwei}

\erg{14}{a) stimmt \quad b) Der Stichprobenanteil liegt stets im Intervall; er ist größer als die Annahme, also ist es erst recht die obere Grenze \quad c) etwa $0{,}04$: vierfacher Umfang, die Länge wird mit $\frac{1}{\sqrt{4}} = \frac{1}{2}$ multipliziert \quad d) wahr für n zwischen 485 und 553: $0{,}28$ ist unverträglich, sobald $n > 484{,}0$ (die Schwelle schrumpft mit n), $0{,}36$ bleibt verträglich, solange $n \le 553{,}2$}
\erg{15}{a) Richtig: etwa $0{,}071$. Fehler: Tim nimmt die Länge proportional zu eins durch n; sie ist proportional zu eins durch Wurzel n \quad b) Setzt man p gleich dem Stichprobenanteil, ist der Abstand null und damit kleiner als jede Schwelle – dieser Wert ist verträglich, liegt also im Intervall}
